% ===== PRÉAMBULE COMMUN =====
\documentclass[11pt,a4paper]{article}
\usepackage[utf8]{inputenc}
\usepackage[T1]{fontenc}
\usepackage{lmodern}
\usepackage[french]{babel}
\usepackage{amsmath,amssymb,mathtools}
\usepackage{icomma}
\usepackage{xcolor}
\usepackage{tikz}
\usepackage{pgfplots}
\usepackage[most]{tcolorbox}
\usepackage{enumitem}
\usepackage{array,booktabs,colortbl}
\usepackage[a4paper,margin=2.2cm,headheight=26pt,headsep=10pt]{geometry}
\usepackage{fancyhdr}
\usepackage[hidelinks]{hyperref}
\usetikzlibrary{arrows.meta,calc,angles,quotes,positioning,decorations.markings,intersections,backgrounds}
\pgfplotsset{compat=1.18}

\definecolor{primary}{RGB}{226,74,88}
\definecolor{primarydark}{RGB}{150,28,42}
\definecolor{methcol}{RGB}{40,90,150}
\definecolor{excol}{RGB}{0,135,90}
\definecolor{defcol}{RGB}{0,114,178}
\definecolor{attncol}{RGB}{200,40,55}
\definecolor{thmcol}{RGB}{0,140,120}
\definecolor{courbe}{RGB}{32,110,200}
\definecolor{courbeder}{RGB}{210,70,40}
\definecolor{tang}{RGB}{0,150,90}
\definecolor{grille}{RGB}{200,205,215}
\definecolor{plus}{RGB}{200,60,60}
\definecolor{moins}{RGB}{30,90,200}

\tcbset{boxbase/.style={enhanced,breakable,boxrule=0.9pt,arc=2.5pt,
  left=3mm,right=3mm,top=2.3mm,bottom=2.3mm,fonttitle=\bfseries,coltitle=white}}
\newtcolorbox{cle}[1][]{boxbase,colback=defcol!6,colframe=defcol,
  title={\textsc{À retenir}\ifx\relax#1\relax\relax\else\textmd{ : \itshape #1}\fi}}
\newtcolorbox{methode}[1][]{boxbase,colback=methcol!7,colframe=methcol,sharp corners,
  title={\textsc{Méthode}\ifx\relax#1\relax\relax\else\textmd{ --- \itshape #1}\fi}}
\newtcolorbox{exemplebox}[1][]{boxbase,colback=excol!5,colframe=excol,
  title={\textsc{Exemple}\ifx\relax#1\relax\relax\else\textmd{ : \itshape #1}\fi}}
\newtcolorbox{piege}[1][]{boxbase,colback=attncol!6,colframe=attncol,
  title={\textsc{Attention}\ifx\relax#1\relax\relax\else\textmd{ : \itshape #1}\fi}}
\newtcolorbox{exo}[1][]{boxbase,colback=primary!4,colframe=primary,
  title={\textsc{Exercice}\ifx\relax#1\relax\relax\else\textmd{~#1}\fi}}
\newtcolorbox{corr}[1][]{boxbase,colback=thmcol!5,colframe=thmcol,
  title={\textsc{Corrigé}\ifx\relax#1\relax\relax\else\textmd{~#1}\fi}}

\newcommand{\R}{\mathbb{R}}
\newcommand{\N}{\mathbb{N}}
\newcommand{\ds}{\displaystyle}
\newcommand{\tg}{\operatorname{tg}}
\newcommand{\cotg}{\operatorname{cotg}}
\newcommand{\arctg}{\operatorname{arctg}}
\newcommand{\arccotg}{\operatorname{arccotg}}
\newcommand{\limh}{\lim_{h\to 0}}

\pagestyle{fancy}\fancyhf{}
\fancyhead[L]{\small\textcolor{primarydark}{\textbf{Thème 5} — Tangente \& lecture graphique}}
\fancyhead[R]{\small\textcolor{primarydark}{\textsc{Corrigé --- Facile}}}
\fancyfoot[C]{\small\thepage}
\renewcommand{\headrulewidth}{0.8pt}

\begin{document}

\begin{center}
{\LARGE\bfseries\color{primarydark} Corrigé --- Écrire l'équation d'une tangente}\\[2pt]
{\color{primary}\rule{0.5\textwidth}{1.2pt}}\\[4pt]
{\small\itshape Niveau Facile \quad$\bullet$\quad 5 exercices}
\end{center}
\vspace{2mm}

\begin{corr}[1 --- $f(x)=x^2$ en $x_0=1$]
\textbf{Point :} $f(1)=1$. \quad\textbf{Pente :} $f'(x)=2x$, donc $f'(1)=2$.\\
\textbf{Tangente :} $y=(x-1)\cdot 2+1=2x-2+1=\boxed{2x-1}$.
\end{corr}

\begin{corr}[2 --- $f(x)=x^3$ en $x_0=2$]
\textbf{Point :} $f(2)=2^3=8$. \quad\textbf{Pente :} $f'(x)=3x^2$, donc $f'(2)=3\cdot 4=12$.\\
\textbf{Tangente :} $y=(x-2)\cdot 12+8=12x-24+8=\boxed{12x-16}$.
\end{corr}

\begin{corr}[3 --- $f(x)=1/x$ en $x_0=1$]
\textbf{Point :} $f(1)=1$. \quad\textbf{Pente :} $f'(x)=-\dfrac{1}{x^2}$, donc $f'(1)=-1$.\\
\textbf{Tangente :} $y=(x-1)\cdot(-1)+1=-x+1+1=\boxed{-x+2}$.
\end{corr}

\begin{corr}[4 --- $f(x)=\sqrt{x}$ en $x_0=4$]
\textbf{Point :} $f(4)=\sqrt{4}=2$. \quad\textbf{Pente :} $f'(x)=\dfrac{1}{2\sqrt{x}}$, donc $f'(4)=\dfrac{1}{2\cdot 2}=\dfrac14$.\\
\textbf{Tangente :} $y=(x-4)\cdot\dfrac14+2=\dfrac{x}{4}-1+2=\boxed{\dfrac{x}{4}+1}$.
\end{corr}

\begin{corr}[5 --- $f(x)=x^2-2x$ en $x_0=0$]
\textbf{Point :} $f(0)=0$. \quad\textbf{Pente :} $f'(x)=2x-2$, donc $f'(0)=-2$.\\
\textbf{Tangente :} $y=(x-0)\cdot(-2)+0=\boxed{-2x}$.
\end{corr}

\end{document}
